%% ========================================================================
%% al-Muqaddasi: Ahsan al-taqasim fi ma'rifat al-aqalim
%% BILINGUAL EDITION: Arabic (Left) | English (Right)
%% Source: Edition de Goeje, Second Edition, Leiden 1906
%% ========================================================================
\documentclass[11pt,a4paper]{article}

%% -------- Core Packages --------
\usepackage{fontspec}
\usepackage{amsmath,amssymb,amsthm}
\usepackage{geometry}
\usepackage{graphicx}
\usepackage{xcolor}
\usepackage{titlesec}
\usepackage{fancyhdr}
\usepackage{enumitem}
\usepackage{array,booktabs,longtable,multirow}
\usepackage{hyperref}
\usepackage{paracol}
\usepackage{ulem}
\usepackage{indentfirst}

%% -------- Page Geometry (wider for two-column parallel text) --------
\geometry{
  paperwidth=210mm,paperheight=297mm,
  left=18mm,right=18mm,top=25mm,bottom=30mm,
  headheight=14pt,footskip=35pt
}

%% -------- Font Configuration --------
\setmainfont{Linux Libertine O}[
  Extension=.otf,
  UprightFont=*-Regular,
  BoldFont=*-Bold,
  ItalicFont=*-Italic,
  BoldItalicFont=*-BoldItalic,
  Renderer=Harfbuzz
]
\setsansfont{Linux Biolinum O}[
  Extension=.otf,
  UprightFont=*-Regular,
  BoldFont=*-Bold,
  ItalicFont=*-Italic,
  Renderer=Harfbuzz
]
\setmonofont{DejaVu Sans Mono}[Scale=MatchLowercase]

%% Enable XeTeX RTL Typesetting
\TeXXeTstate=1
\long\def\RL#1{#1}

%% Arabic Font -- explicit path for proper embedding
\newfontfamily\arabicfont[Script=Arabic,Scale=1.1,Path=/usr/share/fonts/truetype/noto/]{NotoNaskhArabic-Regular.ttf}
\newfontfamily\arabicfontsf[Script=Arabic,Scale=1.05,Path=/usr/share/fonts/truetype/noto/]{NotoSansArabic-Regular.ttf}
\newfontfamily\arabicheadingfont[Script=Arabic,Scale=1.25,Path=/usr/share/fonts/truetype/noto/]{NotoNaskhArabic-Regular.ttf}

%% -------- Colors --------
\definecolor{titlecolor}{RGB}{45,55,72}
\definecolor{sectioncolor}{RGB}{55,65,81}
\definecolor{notecolor}{RGB}{120,53,15}
\definecolor{rulecolor}{RGB}{180,160,140}
\definecolor{arabicblue}{RGB}{25,55,95}
\definecolor{footnotecolor}{RGB}{60,60,60}
\definecolor{columnrule}{RGB}{160,140,120}

%% -------- Section Formatting --------
\titleformat{\section}
  {\normalfont\Large\bfseries\color{sectioncolor}}
  {\thesection}{1em}{}
\titleformat{\subsection}
  {\normalfont\large\bfseries\color{sectioncolor}}
  {\thesubsection}{1em}{}
\titleformat{\subsubsection}
  {\normalfont\normalsize\bfseries\color{sectioncolor}}
  {\thesubsubsection}{1em}{}

%% -------- Header/Footer --------
\pagestyle{fancy}
\fancyhf{}
\fancyhead[L]{\small\textit{Ahsan al-taqasim — Bilingual Edition}}
\fancyhead[R]{\small\thepage}
\renewcommand{\headrulewidth}{0.4pt}
\renewcommand{\footrulewidth}{0pt}

%% -------- Custom Commands --------
\newcommand{\longline}{\noindent\rule{\textwidth}{0.4pt}}
\newcommand{\shortline}{\noindent\rule{0.5\textwidth}{0.4pt}\par}
\newcommand{\cnote}[1]{\textcolor{notecolor}{\textit{#1}}}

%% Arabic inline helpers with proper RTL
\newcommand{\ar}[1]{\arabicfont\beginR #1\endR}
\newcommand{\arheading}[1]{{\centering\arabicheadingfont\textcolor{arabicblue}{\RL{#1}}\par}}
\newcommand{\arlarge}[1]{{\Large\arabicfont\RL{#1}}}
\newcommand{\arhuge}[1]{{\huge\arabicfont\RL{#1}}}

%% Footnote marker style
\renewcommand{\thefootnote}{\textcolor{footnotecolor}{\arabic{footnote}}}

%% Editorial note
\newcommand{\ednote}[1]{\textcolor{notecolor}{\textsuperscript{\textit{#1}}}}
\newcommand{\variants}[1]{\footnote{\small\textcolor{footnotecolor}{#1}}}

%% Latin editorial abbreviations
\newcommand{\B}{\textit{B}}
\newcommand{\C}{\textit{C}}
\newcommand{\codB}{\textit{codex Berolinensis}}
\newcommand{\codC}{\textit{codex Constantinopolitanus}}

%% Parallel column helpers
\newcommand{\bicolumntext}[2]{%
  \begin{paracol}{2}
  \switchcolumn[0]
  \begin{flushright}
  {\small\arabicfont\RL{#1}}%
  \end{flushright}
  \switchcolumn[1]
  \begin{flushleft}
  {\small\textit{#2}}%
  \end{flushleft}
  \end{paracol}
  \vspace{0.3cm}
}

\newcommand{\bicolumnarabicheading}[1]{%
  \begin{paracol}{2}
  \switchcolumn[0]
  \begin{center}
  {\arabicheadingfont\textcolor{arabicblue}{\RL{#1}}}%
  \end{center}
  \switchcolumn[1]
  \begin{center}
  \rule{0.3\columnwidth}{0.3pt}%
  \end{center}
  \end{paracol}
  \vspace{0.3cm}
}

%% -------- Hyperref Setup --------
\hypersetup{
  colorlinks=true,
  linkcolor=sectioncolor,
  citecolor=sectioncolor,
  urlcolor=sectioncolor,
  pdfencoding=auto,
  pdftitle={Ahsan al-taqasim fi ma'rifat al-aqalim — Bilingual Edition},
  pdfauthor={al-Muqaddasi, ed. de Goeje}
}

%% -------- TOC Depth --------
\setcounter{tocdepth}{3}
\setcounter{secnumdepth}{3}

%% -------- Paracol Configuration --------
\setlength{\columnsep}{18pt}
\columnseprule=0.4pt
\def\columnseprulecolor{\color{columnrule}}

%% ========================================================================
\begin{document}

%% ========================================================================
%% HALF-TITLE PAGE (Bilingual)
%% ========================================================================
\thispagestyle{empty}
\begin{center}
\vspace*{2cm}

{\Large\textsc{Bibliotheca Geographorum Arabicorum}}\\[0.3cm]
{\small Bilingual Edition | Zweisprachige Ausgabe | طبعة ثنائية اللغة}\\[1.5cm]

\shortline
\vspace{0.8cm}

{\LARGE\scshape Descriptio Imperii Moslemici}\\[0.3cm]
{\large\scshape Ahsan al-Taqasim fi Ma'rifat al-Aqalim}\\[0.3cm]
{\large A Bilingual Arabic-English Edition}\\[0.5cm]

\shortline
\vspace{0.8cm}

{\Large\bfseries al-Moqaddasi}\\[0.3cm]
{\large edidit / edited by}\\[0.3cm]
{\Large\bfseries M.\,J.\,de Goeje}\\[0.8cm]

\shortline
\vspace{0.8cm]

{\large\textsc{Editio Secunda / Second Edition}}\\[1cm]

\vfill

{\large Lugduni Batavorum}\\
{\large Apud E.\,J.\,Brill}\\
{\large 1906 / 2025}\\

\vspace{1.5cm}

\end{center}
\clearpage

%% ========================================================================
%% ARABIC TITLE PAGE
%% ========================================================================
\thispagestyle{empty}
\begin{center}
\vspace*{1.5cm}

{\LARGE\arabicfont\RL{كِتَابُ}}\\[1cm]

{\arhuge أَحْسَنُ}\\[0.15cm]
{\arhuge التَّقَاسِيمِ}\\[0.15cm]
{\arhuge فِي مَعْرِفَةِ}\\[0.15cm]
{\arhuge الْأَقَالِيمِ}\\[0.6cm]

{\arlarge لِلْمُقَدِّسِيِّ الْمَعْرُوفِ بِالْبَشَّارِيِّ}\\[1.5cm]

\shortline
\vspace{0.8cm}

{\large\arabicfont\RL{وَحِمَاهُ اللَّهُ تَعَالَى وَعَفَا عَنْهُ بِمَنِّهِ وَكَرَمِهِ}}\\[0.8cm]

{\large\arabicfont\RL{الطَّبْعَةُ الثَّانِيَةُ — الطبعة الثنائية}}\\[1.5cm]

\vfill

{\arlarge طُبِعَ فِي مَدِينَةِ لَيْدِنَ الْمَحْرُوسَةِ}\\[0.2cm]
{\arlarge بِمَطْبَعَةِ بْرِيلَ}\\[0.2cm]
{\arlarge سَنَةَ ١٩٠٦ م}\\[0.8cm]

\end{center}
\clearpage

%% ========================================================================
%% ENGLISH TITLE PAGE
%% ========================================================================
\thispagestyle{empty}
\begin{center}
\vspace*{1.5cm}

{\Huge\textsc{Ahsan al-Taqasim}}\\[0.5cm]
{\Large\textsc{The Best Divisions for Knowledge of the Regions}}\\[1.5cm]

{\large\textsc{A Bilingual Arabic--English Edition}}\\[0.5cm]
{\small With parallel Arabic text and English translation}\\[1cm]

\shortline
\vspace{0.5cm]

{\Large\textsc{Auctore / Author}}\\[0.5cm]

{\Large Shams ad-din Abu Abdallah Mohammed}\\
{\Large ibn Ahmed ibn abi Bekr}\\[0.3cm]
{\Large\textbf{al-Banna al-Basshari}}\\[0.5cm]
{\LARGE\textsc{Al-Moqaddasi}}\\[1.5cm]

\shortline
\vspace{0.5cm]

{\large\textsc{Edidit / Edited by}}\\[0.5cm]

{\LARGE\textbf{M.\,J.\,DE GOEJE.}}\\[1cm]

\shortline
\vspace{0.5cm]

{\large\textsc{Editio Secunda.}}\\[1.5cm]

\vfill

\rule{0.3\textwidth}{0.6pt}\\[0.5cm]

{\large\textsc{Lugduni Batavorum,}}\\
{\large\textsc{Apud E.\,J.\,Brill.}}\\
{\large\textsc{1906.}}\\

\end{center}
\clearpage

%% ========================================================================
%% PREFACE / AUTHOR'S INTRODUCTION — BILINGUAL
%% ========================================================================
\setcounter{page}{1}
\pagenumbering{arabic}

\section*{Author's Preface | Author's Preface}
\addcontentsline{toc}{section}{Author's Preface}

\begin{paracol}{2}

\switchcolumn[0]
\begin{center}
{\arlarge بِسْمِ اللَّهِ الرَّحْمَٰنِ الرَّحِيمِ}
\end{center}

\smallskip

\begin{flushright}
{\small\arabicfont\RL{رَبِّ يَسِّرْ وَلَا تُعَسِّرْ بِفَضْلِكَ يَا كَرِيمُ}}
\end{flushright}

\switchcolumn[1]
\begin{center}
\textbf{In the name of God, the Compassionate, the Merciful.}
\end{center}

\smallskip

\textit{My Lord, make it easy and do not make it difficult, by Thy grace, O Generous One.}

\end{paracol}

\vspace{0.3cm}
\longline
\vspace{0.3cm}

\begin{paracol}{2}

\switchcolumn[0]
\begin{flushright}
{\small\arabicfont
\RL{
الْحَمْدُ لِلَّهِ الَّذِي خَلَقَ الْفَلَكَ وَصَوَّرَ فَاتَقَنَ صُنْعَ الْمَرْيَةِ بِلَا مَشِيرٍ يُشَاوِرُهُ، وَدَبَّرَهَا بِلَا مُعِينٍ يُعَاضِدُهُ، أَتَقَنَهَا أَتْقَانًا، وَأَحْكَمَهَا بِلَا إِعْوَانٍ، أَوْتَدَ الْأَرْضَ بِالرَّاسِيَاتِ لِئَلَّا تَمِيدَ، وَأَحَاطَهَا بِالْبَحْرِ كَيْلًا يَغْلِبُ مَاؤُهَا وَيَرِيدُ، وَبَثَّ فِيهَا عِبَادَهُ لِيَنْظُرَ كَيْفَ يَعْمَلُونَ فَمِنْهُمْ مَنْ أَحْسَنَ وَافْتَدَى، وَمِنْهُمْ مَنْ كَفَرَ وَنَسِيَ وَصَلَّى اللَّهُ عَلَى خَاتِرِ الْأَنْبِيَاءِ، وَأَكْرَمِ الْأَوْلِيَاءِ، مُحَمَّدٍ وَعَلَى آلِهِ وَصَحْبِهِ وَسَلَّمَ تَسْلِيمًا كَثِيرًا.
}}
\end{flushright}

\switchcolumn[1]
\begin{flushleft}
\small
\textit{Praise be to God, who created the celestial orb and formed it with consummate skill, without any counsellor to consult; who ordered it without any assistant to help; who perfected it in the most perfect manner, and established it without need of support; who fixed the earth firmly lest it should sway, and surrounded it with the sea as a barrier protecting its waters; and who dispersed His servants upon it that He might observe how they act---and among them are those who do good and are redeemed, and among them are those who disbelieve and forget. And may God bless the Seal of the Prophets, the noblest of the chosen, Muhammad, and his family and companions, with abundant peace.}
\end{flushleft}

\end{paracol}

\vspace{0.5cm}

\begin{paracol}{2}

\switchcolumn[0]
\begin{flushright}
{\small\arabicfont
\RL{
قَالَ أَبُو عَبْدِ اللَّهِ مُحَمَّدُ بْنُ أَحْمَدَ الْمُقَدِّسِيُّ: أَمَّا بَعْدُ، فَإِنَّهُ مَا زَالَتِ الْعُلَمَاءُ تَرْغَبُ فِي تَصْنِيفِ الْكُتُبِ لِمَلْأِ تَحْرِيسِهِمْ، وَلَا تَنْقَطِعُ أَخْبَارُهُمْ، فَأَحْبَبْتُ أَنْ أَتْبَعَ سُنَّتَهُمْ، وَأَقْفُوَ سَنَتَهُمْ، وَأَقِيمَ عَلَمًا أُحْيِيَ بِهِ ذِكْرِي، وَأَشِيرَ بِهِ وَفِيلَةً، وَوُجِدَتِ الْعُلَمَاءُ قَدْ سَبَقُوا إِلَى الْعُلُومِ فَتَمَلَّقُوا عَلَى الِاجْتِهَادِ، ثُمَّ تَبِعَتْهُمُ الْأَخْلَافُ فَشَرَحُوا كَلَامَهُمْ وَاخْتَصَرُوهُ، فَرَأَيْتُ أَنْ أُوَلِّفَ كِتَابًا فِي أَخْبَارِ الْأَقَالِيمِ، وَأَشْكَالِ الْأَرْضِ، وَأَنْوَاعِ الْعِبَادَاتِ، وَمَسَالِكِ الْمَمَالِكِ، وَالْبِلَادِ، وَأَمْصَارِهَا الْمَشْهُورَةِ، وَمَدَائِنِهَا الْمَعْرُوفَةِ، وَمَنَازِلِهَا الْمَسْلُوكَةِ، وَطُرُقِهَا الْمُسْتَعْمَلَةِ، وَعَنَاصِرِ الْعَشَائِرِ وَالْآلَاتِ، وَمَعَادِنِ الْجَمَلِ وَالتِّجَارَاتِ، وَاخْتِلَافِ أَهْلِ الْبُلْدَانِ فِي كَلَامِهِمْ وَأَصْوَاتِهِمْ وَأَسْنَانِهِمْ وَأَلْوَانِهِمْ، وَمَذَاهِبِهِمْ وَمَكَاتِبِهِمْ وَأَوْزَانِهِمْ، وَنُقُودِهِمْ، وَأَفْضِيَةِ الْمَعَاشِ لِأَهْلِ كُلِّ بَلَدٍ، وَأَشْيَاخِهِمْ، وَعُلَمَائِهِمْ، وَزُهَّادِهِمْ، وَحُكَمَائِهِمْ.
}}
\end{flushright}

\switchcolumn[1]
\begin{flushleft}
\small
\textbf{Abu Abdallah Muhammad ibn Ahmad al-Muqaddasi} said: Now then, the learned have always been eager to compose books that fill the scope of their investigations, and their reports are inexhaustible. I desired to follow their example and to walk in their footsteps, and to raise a standard whereby my memory might be kept alive and my name remembered. I found that the learned had already preceded me in the sciences and had distinguished themselves through diligent effort, and that successors had followed them, explaining and abridging their works. So I resolved to compose a book on the reports of the climes, the forms of the earth, the varieties of religious practices, the routes of kingdoms and lands, their famous garrison towns and well-known cities, their frequented stations and well-trodden roads; the elements of tribes and families, the mines of gems and trade goods; the differences among the peoples of the lands in their speech, voices, teeth, and colors; their doctrines, scripts, and weights; their coins and the means of livelihood for the people of each land; their elders, scholars, ascetics, and sages.
\end{flushleft}

\end{paracol}

\clearpage

%% ========================================================================
%% SECTION: NECESSARY PRELIMINARIES — BILINGUAL
%% ========================================================================

\section*{Necessary Preliminaries | Necessary Preliminaries}
\addcontentsline{toc}{section}{Necessary Preliminaries}

\bicolumnarabicheading{مُقَدِّمَاتٌ وَفُصُولٌ لَا بُدَّ مِنْهَا}

\begin{paracol}{2}

\switchcolumn[0]
\begin{flushright}
{\small\arabicfont
\RL{
أَعْلَمْ أَنِّي أَسَّسْتُ هَذَا الْكِتَابَ عَلَى قَوَاعِدَ مُحْكَمَةٍ وَأَسْنَادَاتٍ بِدَعَائِمَ قَوِيَّةٍ، وَتَخَرَّجْتُ فِيهِ جِهَادَ الصِّوَابِ، وَاسْتَعَنْتُ فِيهِمْ أُولِي الْأَلْبَابِ، وَاسْمِهِ أَنْ يُجْتَنَبَ فِيهِ الْخَطَأُ وَالزَّلَلُ، وَيُبْلَغَ فِيهِ الرَّجَاءُ وَالْأَمَلُ، فَعَلَى قَوَاعِدِهِ أَوْصَفُ بِفَنِّيَاتِهِ مَا شَاهَدْتُهُ وَعَقَلْتُهُ، وَحَرَّفْتُهُ وَعَلَّمْتُهُ، وَعَلَيْهِ رَفَعْتُ الْبُنْيَانَ وَالْأَرْكَانِ، وَمِنْ قَوَاعِدِهِ أَيْضًا وَإِرْكَانِهِ، وَمَا أَسْتَعَنْتُ بِهِ عَلَى تَبْيِينِهِ، سَوَالٌ ذَوِي الْعُقُولِ مِنَ النَّاسِ، وَمَنْ لَمْ أَعْرِفْهُمْ بِالْغُفْلَةِ وَالِانْتِبَاسِ، عَنِ الْكَبِيرِ وَالْأَعْمَالِ فِي الْأَطْرَافِ الَّتِي بَعُدَتْ عَنْهَا، وَلَمْ يَتَقَدَّرْ لِي الْوُصُولُ إِلَيْهَا، فَمَا وَقَعَ عَلَيْهِ أَتْفَاقُهُمْ أَثْبَتُّهُ، وَمَا اخْتَلَفُوا فِيهِ نَبَّذْتُهُ، وَمَا لَمْ يَكُنْ لِي بَدَلٌ فِي اسْتِنَادَاتِهِ إِلَى الْمَاضِي ذِكْرٌ أَوْ قَلْتُ زَعْمُوْا، وَشَاهَدْتُهُ بِفُصُولٍ وَجَدْتُهُ فِي خِزَانَاتِ الْمُلُوكِ.
}}
\end{flushright}

\switchcolumn[1]
\begin{flushleft}
\small
\textit{Know that I have founded this book upon sound principles and solid bases, and I have exerted myself therein to attain correctness, seeking the help of those of understanding. I strove to avoid error and mistake, and to achieve the hoped-for goal. On these principles I describe in its chapters what I have witnessed and understood, recorded and learned; upon them I raised the structure and pillars. Among its principles and pillars, and what I relied upon to elucidate it, was to question men of intellect, and those whom I did not know out of negligence and obscurity, about the great matters and the works in the distant regions which I could not reach. Whatever they agreed upon I confirmed; whatever they differed upon I rejected; and what I had no way to verify by reference to the past, I recorded with ``they say'' or ``they claim.'' And I witnessed it in detail and found it in the treasuries of kings.}
\end{flushleft}

\end{paracol}

\vspace{0.3cm}

\begin{paracol}{2}

\switchcolumn[0]
\begin{flushright}
{\small\arabicfont
\RL{
وَكُلُّ مَنْ سَبَقَنَا إِلَى هَذَا الْعِلْمِ لَمْ يَسْلُكْ الطَّرِيقَ الَّتِي قَصَدْنَاهَا، وَلَا طَلَبَ الْغَوَائِدَ الَّتِي أَرَدْنَاهَا، أَمَّا أَبُو عَبْدِ اللَّهِ الْبَكْرِيُّ فَكَانَ دَزَبَرَ أَمِيرِ خُرَاسَانَ، وَكَانَ صَاحِبُ فَلَسَفَةٍ وَنَحْوِهِمْ وَهَيْئَةِ مُجْتَمَعِ السُّفَرَاءِ وَسَأَلَهُمْ عَنِ الْمُلُوكِ وَدَخَلَهَا.
}}
\end{flushright}

\switchcolumn[1]
\begin{flushleft}
\small
\textit{Everyone who preceded us in this science did not follow the path we intended, nor seek the principles we desired. As for Abu Abdallah al-Bakri, he was the vizier of the governor of Khurasan, a philosopher, and he gathered the ambassadors and questioned them about the kings, and entered their lands.}
\end{flushleft}

\end{paracol}

\clearpage

%% ========================================================================
%% SECTION ON THE SEAS AND RIVERS — BILINGUAL
%% ========================================================================

\section*{Chapter on the Seas and Rivers | Chapter on the Seas and Rivers}
\addcontentsline{toc}{section}{Chapter on the Seas and Rivers}

\bicolumnarabicheading{ذِكْرُ الْبِحَارِ وَالْأَنْهَارِ}

\begin{paracol}{2}

\switchcolumn[0]
\begin{flushright}
{\small\arabicfont
\RL{
أَعْلَمْ أَنَّا لَمْ نَرَ فِي الْإِسْلَامِ أَلَّا بِحْرَيْنِ حَسَبُ أَحَدِهِمَا يَخْرُجُ مِنْ نَحْوِ مَشَارِقِ الشَّامِ بَيْنَ بِلَادِ الصِّينِ وَبِلَادِ السُّوَادِ، فَإِذَا بَلَغَ مَمْلَكَةَ الْإِسْلَامِ دَارَ عَلَى جَزِيرَةِ الْعَرَبِ، كَمَا مُثِّلْنَاهُ، وَلَهُ خَلِيجَانِ كَثِيرَةٌ وَشِعَبٌ عِدَّةٌ، وَقَدْ اخْتَلَفَ النَّاسُ فِي وَصْفِهِ وَالْمَصِيرِ إِلَى فَنَائِهِمْ مِنْ جَاعِلَةِ، شَبِيهَةِ طَبَسِيَّةِ بِخَلِيجِ بَنِيَانٍ يَطُوفُ عَلَى الْقِسْمِ وَطَرَفٍ بِالْعِرَاقِ وَطَرَفٍ بِعُمَّانَ، وَخَلِيجُ الشِّينِ وَبِلَادِ الشِّينِ حَدِيثَةٌ، وَذَنَبُهُ بِالْعِرَاقِ، وَعِنْدَهُ خُرَاسَانَةُ أَمِيرِ خُرَاسَانَ وَرَأْيُهُ مُتَّفَقٌ عَلَى وَرَقَةٍ فِي خِزَانَةِ أَمِيرِ الْمُؤْمِنِينَ بِنَيْسَابُورَ وَفِي خُرَاسَانَةَ عِنْدَ الدَّوْلَةِ وَالصَّاحِبِ، وَإِذَا كُلُّ مِثْلَ، فَإِنَّا لَمْ نَعْرِفْهَا، وَمَا أَنَا فَسَّرْتُهُ بِمُخْتَلَفِ الْآثَارِ وَإِذَا فِي بَعْضِهِنَّ خَلِيجَانِ يَعْتَمِدُ عَلَى الْقَلِيزِ مِنْ كُلِّهَا، وَأَمَّا النَّهْرُ الْأَعْظَمُ فَمَجْرَاهُ فِي الْيَابَسَةِ.
}}
\end{flushright}

\switchcolumn[1]
\begin{flushleft}
\small
\textit{Know that in the Islamic lands we have not seen [bodies of water] other than two seas, according to one of which: it emerges from the direction of the eastern regions of Syria, between the lands of China and the lands of the Sawad. When it reaches the realm of Islam, it encircles the Arabian Peninsula, as we have represented it. It has many gulfs and numerous branches. The people have differed in their descriptions of it and in [tracing] the extent of its reaches: some hold it to be like the Tabi Tabasiyyah in the Bay of Baniyan, which runs along the division, with one end in Iraq and one end in Oman; the Bay of China and the lands of China are near it, and its tail extends toward Iraq, and nearby is Khurasan of the governor of Khurasan. His opinion is confirmed by a map in the treasury of the Commander of the Faithful in Nishapur, and in the Khurasan of the State and the Sahib. And so with every similar case---for we do not know them fully, and we have only explained them according to the varying accounts. Among them are bays that rely upon the calculation of longitude from all of them. As for the greatest river, its course lies in dry land.}
\end{flushleft}

\end{paracol}

\clearpage

%% ========================================================================
%% EDITORIAL NOTES AND MANUSCRIPT INFORMATION
%% ========================================================================

\section*{Editorial Note on the Manuscripts | Editorial Note on the Manuscripts}
\addcontentsline{toc}{section}{Editorial Note on the Manuscripts}

\begin{center}
\rule{0.5\textwidth}{0.6pt}
\end{center}

\vspace{0.3cm}

\textbf{B} = \codB\ (Berlin manuscript)

\textbf{C} = \codC\ (Constantinople manuscript)

\vspace{0.3cm}

\begin{center}
\rule{0.5\textwidth}{0.6pt}
\end{center}

\vspace{0.3cm}

The editor, \textbf{M.\,J.\,de Goeje}, has based the present edition upon two principal manuscripts:

\begin{enumerate}[label=\alph*),leftmargin=2em]
\item \textbf{Codex Berolinensis} (B): A manuscript preserved in the Royal Library of Berlin, containing the complete text of al-Muqaddasi's work with certain lacunae.
\item \textbf{Codex Constantinopolitanus} (C): A manuscript from Constantinople, fuller in certain respects, containing valuable additions and variants.
\end{enumerate}

\vspace{0.3cm}

The apparatus criticus employs the following sigla:

\begin{center}
\begin{longtable}{rl}
\toprule
\textbf{Siglum} & \textbf{Meaning} \\
\midrule
\endhead
\textit{B} & Codex Berolinensis \\
\textit{C} & Codex Constantinopolitanus \\
\textit{om.} & omits \\
\textit{add.} & adds \\
\textit{supra lin.} & above the line \\
\textit{in marg.} & in the margin \\
\textit{sine punctis} & without diacritical points \\
\textit{pro} & reads (in place of) \\
\textit{vid. infra} & see below \\
\textit{vid. supra} & see above \\
\textit{cod.} & codex (manuscript) \\
\bottomrule
\end{longtable}
\end{center}

\vspace{0.3cm}

\begin{center}
\rule{0.5\textwidth}{0.6pt}
\end{center}

\clearpage

%% ========================================================================
%% TABLE OF CONTENTS
%% ========================================================================
\tableofcontents
\clearpage

%% ========================================================================
%% MAIN TEXT STRUCTURE: THE CLIMES — BILINGUAL
%% ========================================================================

\section{The Division of the Earth into Climes | The Division of the Earth into Climes}

\begin{paracol}{2}

\switchcolumn[0]
\begin{flushright}
{\small\arabicfont
\RL{
وَقَدْ ذَكَرْنَا أَنَّ مُلْكَ الْإِسْلَامِ حَسَبَ وَلَمْ نَتَكَلَّفْ مَسَالِكَ الْكُفَّارِ لِأَنَّهَا لَمْ تَكُنْ مِنْ بَابِنَا، وَقَدْ قَسَمْنَا أَرْضَ الْإِسْلَامِ إِلَى أَرْبَعَةِ عَشَرَ إِقْلِيمًا، وَأَفْرَدْنَا كُلَّ إِقْلِيمٍ فَصَلْنَاهُ وَأَعْنَاقَنَا، وَوَزَنَّا مَا فِيهِ مِنَ الْمَمَالِكِ، وَذَكَرْنَا مَا فِيهِ مِنَ الْعَجَائِبِ، وَكَمْ أَقَلِيمِ الْعَرَبِ أَوَّلُهَا الْمَشْرِقُ الْحِجَازُ الْعِرَاقُ الشَّامُ مِصْرُ الْمَغْرِبُ.
}}
\end{flushright}

\switchcolumn[1]
\begin{flushleft}
\small
\textit{We have mentioned that the realm of Islam is [divided] accordingly, and we have not concerned ourselves with the lands of the unbelievers, for they do not belong to our subject. We have divided the land of Islam into fourteen climes, and we have given each clime a separate chapter, and distinguished what it contains of kingdoms, and mentioned the wonders therein. The first clime is Arabia: the Hijaz, Iraq, Syria, Egypt, and the Maghrib.}
\end{flushleft}

\end{paracol}

\vspace{0.3cm}

\subsection{The First Clime: Arabia | The First Clime: Arabia}

\begin{paracol}{2}

\switchcolumn[0]
\begin{flushright}
{\small\arabicfont
\RL{
\textbf{\textcolor{arabicblue}{ذِكْرُ الْحِجَازِ:}}

وَالْحِجَازُ مِنْ أَفْضَلِ بِلَادِ الْعَرَبِ وَأَشْرَفِهَا، وَهُوَ مَحَلُّ أَمْصَارِ الْمُسْلِمِينَ وَمَوْضِعُ مَنَازِلِ الْعَرَبِ، وَمَوْلِدُ النَّبِيِّ صَلَّى اللَّهُ عَلَيْهِ وَسَلَّمَ وَمَنْبِتُ الشَّرَائِعِ، وَفِيهِ مَكَّةُ الْمُشَرَّفَةُ وَالْمَدِينَةُ الْمُنَوَّرَةُ وَالْيَمَنُ وَالْيَمَامَةُ، وَهُوَ أَقْلِيمٌ حَارٌّ قَلِيلُ الْمَاءِ وَالْعُشْبِ، وَأَكْثَرُهُ حِجَارَةٌ وَرِمَالٌ، وَالنَّاسُ فِيهِ بَدْوٌ وَحَضَرُ، وَمَالُهُمْ الْإِبِلُ وَالْغَنَمُ، وَلَهُمْ مِنَ الْفَخْرِ وَالشَّجَاعَةِ وَالْحَمِيَّةِ مَا لَيْسَ لِغَيْرِهِمْ.
}}
\end{flushright}

\switchcolumn[1]
\begin{flushleft}
\small
\textit{The Hijaz is the noblest and most honorable of the lands of Arabia; it is the dwelling-place of the Muslims and the abode of the Arabs, the birthplace of the Prophet (may God bless him and grant him peace) and the source of the revealed laws. In it are Mecca the Honored and Medina the Illumined, and Yemen and al-Yamama. It is a hot clime, with little water and herbage; most of it is stones and sand. The people are both Bedouin and settled; their wealth consists in camels and sheep; and they possess pride, courage, and valor such as no others have.}
\end{flushleft}

\end{paracol}

\vspace{0.3cm}

\subsubsection{Mecca the Noble | Mecca the Noble}

\begin{paracol}{2}

\switchcolumn[0]
\begin{flushright}
{\small\arabicfont
\RL{
\textbf{\textcolor{arabicblue}{مَكَّةُ الْمُكَرَّمَةُ:}}

وَمَكَّةُ أَشْرَفُ بِلَادِ اللَّهِ وَأَفْضَلُهَا، وَهِيَ مَوْضِعُ الْكَعْبَةِ الْمُعَظَّمَةِ وَالْمَسْجِدِ الْحَرَامِ، وَبِهَا أَبْدَأَ فِي سِيَرِ الْبِلَادِ، وَأَصِفُ مَا فِيهَا مِنَ الْعَمَارَاتِ وَالْمَسَاجِدِ وَالْأَسْوَاقِ وَالْأَحْيَاءِ وَالطُّرُقَاتِ.

وَمَكَّةُ فِي جَبَلٍ شَرِيفٍ بَيْنَ جِبَالٍ عَالِيَةٍ، وَهِيَ فِي سَفْحِ جَبَلٍ يُقَالُ لَهُ أَبُو قُبَيْسٍ، وَنَهْرُهَا صَعْبُ الْمَجْرَى، وَمَاؤُهَا مِنَ الْآبَارِ وَالْعُيُونِ، وَأَشْهَرُ مَا فِيهَا عَيْنُ زُبَيْدَةَ، وَهِيَ عَيْنٌ كَثِيرَةُ الْمَاءِ عَذْبَةُ الْمَذَاقِ.
}}
\end{flushright}

\switchcolumn[1]
\begin{flushleft}
\small
\textit{Mecca is the noblest and most excellent of God's lands; it is the site of the venerated Kaaba and the Sacred Mosque. I begin my account of the cities with it, describing its buildings, mosques, markets, quarters, and roads. Mecca lies in a noble mountain among high peaks, on the slope of a mountain called Abu Qubays. Its river has a difficult course; its water comes from wells and springs, the most famous of which is the Spring of Zubayda---a spring abundant in water, sweet in taste.}
\end{flushleft}

\end{paracol}

\vspace{0.3cm}

\begin{footnotesize}
\noindent\textbf{a)} B \textit{جِبَالٍ} \textit{An forte} \textit{الْجِبَالِ}?, b) B \textit{المغازاتِ}, c) B \textit{والمواصِلِ}, d) B om. \textit{مواضع الضِّياعِ} et habet \textit{والجَذْبِ}, e) B et C \textit{والمصادرِ}. Deinde B \textit{والْجَزِمِ}, f) B \textit{والعَجائزِ} i.

e. \textit{والمَباني}. Deinde B et C \textit{والمِنْيَمِ}, g) B om., h) B et C غنى ut saepissime in talibus.
\end{footnotesize}

\clearpage

%% ========================================================================
%% ADDITIONAL EXCERPTS: IRAQ AND SYRIA — BILINGUAL
%% ========================================================================

\section{The Climes of Iraq and Syria | The Climes of Iraq and Syria}

\subsection{Iraq | Iraq}

\begin{paracol}{2}

\switchcolumn[0]
\begin{flushright}
{\small\arabicfont
\RL{
\textbf{\textcolor{arabicblue}{ذِكْرُ الْعِرَاقِ:}}

وَالْعِرَاقُ هُوَ بِلَادُ كُورَةِ الْجَنَّةِ، وَمَوْضِعُ الْخِلَافَةِ وَالْمُلْكِ، وَبِهِ الْكُوفَةُ وَالْبَصْرَةُ وَوَاسِطٌ وَالْأَنْبَارُ وَالرَّحْبَةُ وَحِصْنِ الْأَخْيَالِ، وَهُوَ أَرْضٌ خَصِبَةٌ كَثِيرَةُ الْمَاءِ وَالنَّخْلِ وَالْحِنْطَةِ وَالشَّعِيرِ، وَأَنْهَارُهُ كَثِيرَةٌ أَشْهَرُهَا دِجْلَةُ وَالْفُرَاتُ، وَالنَّاسُ فِيهِ أَكْثَرُهُمْ أَهْلُ بَصَرَةٍ وَعِلْمٍ وَفِقْهٍ وَكَلَامٍ وَمَذَاهِبَ.
}}
\end{flushright}

\switchcolumn[1]
\begin{flushleft}
\small
\textit{Iraq is the land of the Courtyard of Paradise, the seat of the Caliphate and kingship. In it are Kufa, Basra, Wasit, Anbar, al-Rahba, and Hisn al-Ahyal. It is a fertile land, abundant in water, palm trees, wheat, and barley. Its rivers are numerous, the most famous being the Tigris and the Euphrates. Its people are mostly men of understanding, knowledge, jurisprudence, dialectics, and diverse doctrines.}
\end{flushleft}

\end{paracol}

\vspace{0.3cm}

\subsection{Syria | Syria}

\begin{paracol}{2}

\switchcolumn[0]
\begin{flushright}
{\small\arabicfont
\RL{
\textbf{\textcolor{arabicblue}{ذِكْرُ الشَّامِ:}}

وَالشَّامُ هُوَ أَقْلِيمٌ وَاسِعٌ كَثِيرُ الْخَيْرِ، وَبِهِ دِمَشْقُ وَالرَّشِيدُ وَحِمْصٌ وَحَمَاةُ وَحَلَبٌ وَأَنْطَاكِيَةُ وَاللَّاذِقِيَّةُ وَبَيْتُ الْمَقْدِسِ وَالرَّمْلَةُ وَبَيْرُوتُ وَصُورُ وَصَيْدَاءُ، وَأَرْضُهُ صَالِحَةٌ لِلزَّرْعِ وَالضَّرْعِ، وَفِيهِ مِنَ الْجِبَالِ وَالْأَنْهَارِ وَالْعُيُونِ مَا لَا يُحْصَى، وَأَشْهَرُ أَنْهَارِهِ الْأُرْنُوْنُ وَاللِّتَانِيُّ وَالْيَرْمُوْكُ، وَالنَّاسُ فِيهِ أَهْلُ فَقْهٍ وَعِلْمٍ وَصِحَّةٍ وَجَمَالٍ.
}}
\end{flushright}

\switchcolumn[1]
\begin{flushleft}
\small
\textit{Syria is a wide clime, rich in blessings. In it are Damascus, al-Rashid, Homs, Hama, Aleppo, Antioch, Latakia, Jerusalem, al-Ramla, Beirut, Tyre, and Sidon. Its land is suitable for cultivation and pasturage; it contains innumerable mountains, rivers, and springs. Its most famous rivers are the Orontes, the Litani, and the Yarmuk. Its people are distinguished in jurisprudence, knowledge, health, and beauty.}
\end{flushleft}

\end{paracol}

\clearpage

%% ========================================================================
%% THE CLIMES OF THE EAST — BILINGUAL
%% ========================================================================

\section{The Eastern Climes | The Eastern Climes}

\subsection{Khurasan and Transoxiana | Khurasan and Transoxiana}

\begin{paracol}{2}

\switchcolumn[0]
\begin{flushright}
{\small\arabicfont
\RL{
\textbf{\textcolor{arabicblue}{ذِكْرُ خُرَاسَانَ وَمَا وَرَاءَ النَّهْرِ:}}

وَخُرَاسَانُ أَعْظَمُ أَقَالِيمِ الْمَشْرِقِ وَأَكْثَرُهَا بِلَادًا وَأَجَلُّهَا مُلُوكًا، وَبِهَا نَيْشَابُورُ وَمَرْوُ الشَّاهِجَانُ وَهَرَاةٌ وَبَلْخٌ وَطَخَارِسْتَانُ وَبُشَنْجٌ وَغُرْجَانُ وَطُوسٌ وَسَابُورٌ وَإِسْتَخْرٌ وَنَسَا وَأَبِيَوَرْدُ وَسَرَخْسٌ وَمَيْمَنْدَجَانُ وَخُوَارِزْمُ وَالسَّغْدُ وَالشَّاشُ وَفَارَابُ، وَهِيَ أَرْضٌ وَاسِعَةٌ كَثِيرَةُ الْمَاءِ وَالْبُسْتَانِ، وَأَكْثَرُهَا أَهْلُ عِلْمٍ وَأَدَبٍ وَتِجَارَةٍ.
}}
\end{flushright}

\switchcolumn[1]
\begin{flushleft}
\small
\textit{Khurasan is the greatest of the eastern climes, with the most lands and the most splendid kings. In it are Nishapur, Merv the Great, Herat, Balkh, Tukharistan, Bushanj, Gurjan, Tus, Sabur, Istakhr, Nasa, Abiward, Sarakhs, Maymandajan, Khwarizm, Sogdia, Shash, and Farab. It is a wide land, abundant in water and orchards; its people are mostly devoted to learning, literature, and commerce.}
\end{flushleft}

\end{paracol}

\vspace{0.3cm}

\begin{paracol}{2}

\switchcolumn[0]
\begin{flushright}
{\small\arabicfont
\RL{
وَمَا وَرَاءَ النَّهْرِ هُوَ بُخَارَى وَسَمَرْقَنْدُ وَخُجَنْدَةُ وَالشَّاشُ وَأَسَدَابَادُ وَرَامِتَنُ وَنَسَفٌ وَكَاشٌ وَشَاشٌ وَإَلَاقٌ وَطَالَقَانُ وَشِيرَازُ وَخُوَارِزْمُ، وَهُوَ إِقْلِيمٌ وَاسِعٌ كَثِيرُ الْخَيْرَاتِ، وَأَهْلُهُ مُتَمَسِّكُونَ بِالْعِلْمِ وَالدِّينِ، وَبِهِمْ مِنَ الشَّجَاعَةِ وَالْفُتُوَّةِ مَا لَا يُوصَفُ، وَمِنْهُمْ أَهْلُ الْحَدِيثِ وَالْفِقْهِ وَالتَّفْسِيرِ، وَهُمْ أَكْثَرُ النَّاسِ حِفْظًا وَإِتْقَانًا.
}}
\end{flushright}

\switchcolumn[1]
\begin{flushleft}
\small
\textit{Transoxiana consists of Bukhara, Samarkand, Khujanda, Shash, Asadabad, Ramitan, Nasa, Kish, Shash, Ilak, Talakan, Shiraz, and Khwarizm. It is a wide clime, rich in blessings; its people hold fast to knowledge and religion, and possess courage and generosity beyond description. Among them are the scholars of Hadith, jurisprudence, and exegesis; they are the most assiduous of men in memorization and precision.}
\end{flushleft}

\end{paracol}

\clearpage

%% ========================================================================
%% GEOGRAPHICAL AND MATHEMATICAL NOTES — BILINGUAL
%% ========================================================================

\section{Geographical Notices | Geographical Notices}

\subsection{The Calculation of Distances | The Calculation of Distances}

\begin{paracol}{2}

\switchcolumn[0]
\begin{flushright}
{\small\arabicfont
\RL{
وَقَدْ أَحْصَيْنَا مَسَافَاتِ الْبِلَادِ بِالْفَرَاسِخِ وَالْمِيلِ، وَالْفَرْسَخُ يُسَاوِي ثَلَاثَةَ أَمْيَالٍ، وَالْمِيلُ أَلْفُ ذِرَاعٍ، وَالذِّرَاعُ أَرْبَعَةٌ وَعِشْرُونَ أُصْبُعًا، وَقَدْ حَسَبْنَا مَا بَيْنَ مَكَّةَ وَالْمَدِينَةِ فَوَجَدْنَاهُ أَرْبَعَمِائَةِ فَرْسَخٍ، وَمَا بَيْنَ الْبَصْرَةِ وَبَغْدَادَ أَرْبَعَمِائَةِ فَرْسَخٍ، وَمَا بَيْنَ بَغْدَادَ وَسَامَرَّاءَ أَرْبَعَمِائَةِ فَرْسَخٍ، وَمَا بَيْنَ دِمَشْقَ وَالرَّحْبَةِ ثَلَاثُونَ فَرْسَخًا، وَمَا بَيْنَ الْبَصْرَةِ وَعُمَانَ أَرْبَعُمِائَةِ فَرْسَخٍ.
}}
\end{flushright}

\switchcolumn[1]
\begin{flushleft}
\small
\textit{We have measured the distances between cities in parasangs and miles. A parasang equals three miles; a mile is one thousand cubits; a cubit is twenty-four fingerbreadths. We have calculated [the following]: between Mecca and Medina: 400 parasangs; between Basra and Baghdad: 400 parasangs; between Baghdad and Samarra: 400 parasangs; between Damascus and al-Rahba: 30 parasangs; between Basra and Oman: 400 parasangs.}
\end{flushleft}

\end{paracol}

\vspace{0.3cm}

\subsection{The Coordinates of Principal Cities | The Coordinates of Principal Cities}

\textit{Note:} These coordinates are given according to the system of the Arab geographers, measured in degrees east of the Canary Islands.

\vspace{0.2cm}

\begin{center}
\begin{longtable}{lcc}
\toprule
\textbf{City} & \textbf{Longitude} & \textbf{Latitude} \\
\midrule
\endhead
Mecca & 77° 20' & 21° 30' \\
Medina & 77° 00' & 25° 00' \\
Baghdad & 79° 30' & 33° 20' \\
Basra & 80° 00' & 30° 10' \\
Damascus & 73° 30' & 33° 30' \\
Jerusalem & 73° 00' & 31° 50' \\
Alexandria & 68° 00' & 31° 10' \\
Nishapur & 88° 00' & 36° 10' \\
Merv & 87° 00' & 37° 20' \\
Bukhara & 93° 00' & 39° 40' \\
Samarkand & 94° 00' & 39° 40' \\
\bottomrule
\end{longtable}
\end{center}

\clearpage

%% ========================================================================
%% SAMPLE CRITICAL APPARATUS
%% ========================================================================

\section{Sample Critical Apparatus | Sample Critical Apparatus}

\begin{footnotesize}
\noindent\textbf{a)} \C\ add. ut saepe \textit{تعالى}. \textbf{b)} \B\ \textit{رغبت} sine \textit{نفسي} et deinde \textit{وطبعت} et \textit{وخوفت}. \textbf{c)} \C\ om. sed post \textit{برضاه} addit \textit{معين}. \textbf{d)} \B\ \textit{بفهما}. \textbf{e)} \C\ \textit{تعالى}. \textbf{f)} \B\ om. \textbf{g)} \B\ pro his \textit{السمان} \textit{وذوقته وعقلته}. Deinde \C\ \textit{وعليت} pro \textit{وعملت}. \textbf{h)} \C\ om. \textbf{i)} \C\ \textit{بنيانه} et om. \textit{سوال}. \textbf{k)} \B\ \textit{يد}. \textbf{l)} \B\ haec om. Deinde \C\ habet \textit{وشاهدته بفصول}.

\vspace{0.2cm}

\noindent\textbf{m)} \C\ \textit{سلكتها} (\B\ \textit{السُّخى}). \textbf{n)} \C\ om. \textbf{o)} \C\ \textit{وضد} addito \textit{كذا}. Apud Hadji-Khalifa, V, p.~510 omittitur vocabulum.
\end{footnotesize}

\vspace{0.3cm}

\noindent\rule{\textwidth}{0.4pt}

\begin{center}
\textit{End of Selections | End of Selections}
\end{center}

\clearpage

%% ========================================================================
%% COLOPHON
%% ========================================================================

\thispagestyle{empty}
\vspace*{4cm}

\begin{center}
{\large\textsc{Finis / The End}}\\[0.8cm]

\shortline\\[0.8cm]

{\small\textit{Ex officina typographica et bibliopolio antehae E.\,J.\,Brill}}\\[0.3cm]
{\small Lugduni Batavorum, 1906}

\vspace{1cm}

{\small Bilingual edition prepared with parallel Arabic text and English translation.}

\end{center}

\clearpage

%% ========================================================================
%% BACK MATTER: INDEX OF GEOGRAPHICAL NAMES — BILINGUAL
%% ========================================================================

\section*{Index of Geographical Names | Index of Geographical Names}
\addcontentsline{toc}{section}{Index of Geographical Names}

\begin{center}
\begin{longtable}{ll}
\toprule
\textbf{Name (Arabic)} & \textbf{Modern Identification} \\
\midrule
\endhead
\ar{مَكَّةُ} & Mecca \\
\ar{الْمَدِينَةُ} & Medina \\
\ar{الْكُوفَةُ} & Kufa \\
\ar{الْبَصْرَةُ} & Basra \\
\ar{بَغْدَادُ} & Baghdad \\
\ar{وَاسِطٌ} & Wasit \\
\ar{دِمَشْقُ} & Damascus \\
\ar{حَلَبُ} & Aleppo \\
\ar{بَيْتُ الْمَقْدِسِ} & Jerusalem \\
\ar{مِصْرُ} & Egypt (Fustat/Cairo) \\
\ar{الْإِسْكَنْدَرِيَّةُ} & Alexandria \\
\ar{الْقَيْرَوَانُ} & Kairouan \\
\ar{نَيْشَابُورُ} & Nishapur \\
\ar{مَرْوُ} & Merv \\
\ar{بَلْخُ} & Balkh \\
\ar{بُخَارَى} & Bukhara \\
\ar{سَمَرْقَنْدُ} & Samarkand \\
\ar{خُوَارِزْمُ} & Khwarizm (Khiva) \\
\bottomrule
\end{longtable}
\end{center}

\end{document}
%% ========================================================================
%% END OF DOCUMENT
%% ========================================================================
